\documentclass[a4paper]{article}
\usepackage[a4paper, top=8mm, bottom=8mm, left=8mm, right=8mm]{geometry}
\usepackage{polyglossia}
\setdefaultlanguage[babelshorthands=true]{russian}
\setotherlanguage{english}
\usepackage{fontspec}
\setmainfont{FreeSerif}
\usepackage{parskip}

\title{Отчёт о выездном дне в офисе YADRO}
\author{Белимов Максим Андреевич}
\date{}

\begin{document}

\maketitle

Группа: 25.Б71-мм \hfill Кафедра: системного программирования \hfill Руководитель: Смирнов К.К. \hfill Семестр: 2

\vspace{0.2cm}

21 апреля 2026 года в рамках учебной практики состоялся выездной день в офис компании YADRO.

\textbf{Основные направления.} В первой лекции рассказали о разработке систем хранения данных (СХД) семейства TATLIN, серверного оборудования VEGMAN и роли YADRO в импортозамещении. Лектор объяснил, как разделение данных повышает надёжность и производительность.

\textbf{Стажировка «Импульс».} Вторая лекция была посвящена летней стажировке: требования к кандидатам, отбор (тестовое задание, собеседование), содержание программы. Стажировка ориентирована на студентов со 2 курса.

\textbf{Жизненный цикл разработки.} Третья лекция — о жизненном цикле продукта на примере TATLIN BACKUP: от требований до сопровождения, включая многоуровневое тестирование (от модульных тестов до системного QA). В моей курсовой работе (интерпретатор miniML) используются те же принципы тестирования (Alcotest, QCheck).

\textbf{Q\&A-сессия.} В завершение встретились с сотрудниками — выпускниками СПбГУ. Они рассказали о своём пути в компанию, реальных задачах и совмещении работы с учёбой.

\textbf{Заключение.} Выездной день оставил положительные впечатления, помог лучше представить работу в IT-компании. Стажировка «Импульс» стала реальным ориентиром. Благодарим сотрудников YADRO за открытость.

\end{document}
